\documentclass[11pt]{article}
\usepackage[a4paper,margin=25mm]{geometry}
\usepackage{amsmath,amssymb,amsthm,hyperref}
\newtheorem{theorem}{Theorem}\newtheorem{lemma}{Lemma}
\title{Any fixed number of simultaneous polynomial-size dice}
\author{Complete proof; three independent reviews passed}\date{30 September 2026}
\begin{document}\maketitle
\noindent\textbf{Status.} The full global induction and every input below
have passed independent mathematical reviews by the root, size-bound,
and rational-design agents. This is not formal verification or a
literature-priority claim.

\begin{theorem}\label{main}
For every fixed positive integer $k$ there are constants $C_k,c_k>0$
such that, for all sufficiently large $n$, a fully permutation-fair
family of $n$ finite equiprobable dice exists with $k$ specified dice
having at most $C_k n^{c_k}$ faces each. The remaining $n-k$ dice may
have one common finite face count. No bound on that common count is
asserted, and no uniform control in growing $k$ is asserted.
\end{theorem}

Throughout the proof, polynomial bounds are in $n$, with constants and
exponents allowed to depend on the fixed $k$. In particular a finite
number of polynomial constructions remains polynomial; this convention
does not permit taking the least common multiple of independently
produced polynomial-sized denominators.

\begin{lemma}[One insertion with a reserve of low moments]\label{insert}
Fix $s\ge1$. Suppose $j$ rational probability densities $f_1,\ldots,f_j$
are supported on pure cells of one uniform rational grid of mesh $1/Q$:
every open grid cell is occupied by exactly one density. Suppose:
\begin{enumerate}
\item $Q$, all polynomial piece degrees, and the inverse of every
positive cell mass are polynomially bounded;
\item every normalized conditional density, rescaled to $[0,1]$, lies
between fixed positive constants;
\item all raw unnormalized cell moments through degree $2s+1$ have one
polynomial common denominator $D$.
\end{enumerate}
Then there are $j+1$ rational densities on a refined collection of pure
rational cells, with the same bounds, whose ordered polynomial moments
through any prescribed polynomial degree $M$ coincide with those of
independent variables having densities $f_1,\ldots,f_j,1$.
All new cell endpoints and all new raw cell moments through degree $s$
have one polynomial common denominator. Applying uniform-grid refinement
restores the initial uniform-cell form without changing any of these
ordered moments or the polynomial bounds.
\end{lemma}
\begin{proof}
For an old cell $I=[A/Q,(A+1)/Q]$, let $\gamma$ be its positive mass
and $w_I$ its normalized conditional density in the coordinate
$t=Qx-A$. Its moments satisfy
\begin{equation}\label{normalization}
 \mu_{I,a}:=\int_0^1t^aw_I(t)dt
 =\gamma^{-1}\sum_{l=0}^a\binom al Q^l(-A)^{a-l}
                              \int_I x^lf_i(x)dx.
\end{equation}
Choose one polynomial integer $N$, divisible by $D$, sufficiently large
that all integers $K_I=N\gamma$ satisfy the count threshold of the
synchronized local completion theorem. This is possible because the
cell masses are at least $1/D$ and all local complexity bounds are
polynomial. For $a\le2s+1$, equation~\eqref{normalization} gives
\begin{equation}\label{integer}
 K_I\mu_{I,a}
 =N\sum_{l=0}^a\binom al Q^l(-A)^{a-l}\int_I x^lf_i(x)dx
 \in\mathbb Z.
\end{equation}
Thus the low-moment arithmetic hypothesis of
\path{../size_bounds/synchronized_local_moment_completion.tex}
holds simultaneously for the entire polynomial-sized family of local
inputs. The degrees and conditional density bounds also satisfy its
hypotheses uniformly.

Apply that theorem with output reserve $s$ and ordered exactness $M$.
It supplies one common polynomial denominator $L$ for all low gap
moments of the new uniform variable, one common polynomial denominator
$E$ for all low \emph{conditional} moments of the old-variable boxes,
and one common polynomial endpoint grid. These common denominators
are shared by every local instance; they do not result from taking
least common multiples after independent constructions. Each old box
has normalized mass $1/K_I$.

Rescale the outputs back to $I$. The old variable's global box mass
is $\gamma/K_I=1/N$, and its raw moment of degree $a\le s$ is
\[
 \frac1{NQ^a}\sum_{l=0}^a\binom al A^{a-l}\nu_l,
\]
where $\nu_l$ is the normalized box moment with denominator dividing
$E$. All such moments therefore have common denominator dividing
$NQ^sE$. The new variable's global mass on the macrocell is $1/Q$;
its raw moment on a new local gap is
\[
 \frac1{Q^{a+1}}\sum_{l=0}^a\binom al A^{a-l}\rho_l,
\]
where the local unnormalized gap moments $\rho_l$ have denominator
dividing $L$. Their common denominator divides $Q^{s+1}L$.
All global endpoints have the common local grid denominator multiplied
by $Q$. Multiplying this fixed number of common denominators gives a
single polynomial denominator for the new global cells and their low
moments. The quantitative local construction gives inverse-polynomial
cell lengths, controlled conditional density bounds, and polynomial
piece degrees. Positive masses are at least the reciprocal of the
common denominator.

Assemble the old outputs by their existing labels and the new uniform
outputs as label $j+1$. The resulting densities have total mass one.
To check ordered moments, partition their product domain according to
original grid cells. The $j$ old variables occupy distinct cells because
the cells were pure. The new variable can therefore share a cell with
at most one old variable. If it shares none, only marginal polynomial
moments are needed; if it shares one, the two-variable ordered local
moment identity applies in that cell. All relative orders between
other cells are fixed. Expand a degree-$M$ test polynomial into its
monomials and apply these local identities and the unchanged cell
masses. Summing proves every required global ordered moment identity.

Finally apply
\path{../size_bounds/uniform_grid_moment_refinement.tex}
with low reserve $s$ and high degree $M$. Its input assumptions hold:
there is one polynomial endpoint grid, a common polynomial denominator
for the low moments, inverse-polynomial cell lengths, and polynomial
upper and reciprocal lower density bounds. The latter follow from the
conditional bounds and the polynomial mass/length bounds. The refinement
preserves every original cell moment through degree $M$, retains pure
labels and the low common denominator, and makes the ambient cells
uniform. Thus it preserves all ordered degree-$M$ moments and finishes
the lemma.
\end{proof}

\begin{proof}[Proof of Theorem~\ref{main}]
Fix $k$. Put
\[
 q_j=2^{k-j+1}-1\quad(1\le j\le k),
 \qquad q_j=2q_{j+1}+1.
\]
These are fixed integers, with $q_k=1$. Work throughout with ordered
polynomial exactness $M=n+q_1$; using this larger degree is harmless.
Begin with the single density $f_1=1$ on the one-cell grid. It has all
required conditional bounds, and its moments through $q_1$ have a
fixed common denominator. Together with $n-1$ independent uniforms it
is fully permutation-fair.

Inductively, after constructing $j$ special densities, suppose their
pure cells lie on a uniform polynomial grid and their raw moments
through degree $q_j$ have one polynomial common denominator. Apply
Lemma~\ref{insert} with $s=q_{j+1}$ to separate one of the remaining
uniforms as the next special density. The reserve identity gives exactly
the needed input moments. The lemma preserves every ordered polynomial
moment through degree $M$ of the $j+1$ special variables and restores
the same invariant with reserve $q_{j+1}$. In a fixed full order of
all $n$ variables, integration over the remaining $n-j-1$ uniforms
leaves, on each special-variable chamber, a polynomial of degree
$n-j-1\le M$. Hence full permutation fairness is preserved.
There are only $k-1$ insertions, so every complexity and common
denominator bound remains polynomial in $n$, with exponents depending
on $k$.

After the last step there are $k$ rational pure densities fully fair
with $n-k$ uniforms. Their cell masses have a common polynomial
denominator $D_*$. Choose a common polynomial factor $H$ so that on
each cell of mass $a/D_*$ there is a real equal-weight quadrature with
exactly $Ha$ distinct interior nodes, exact through the larger of degree
$n$ and one more than the polynomial degree of that cell density.
The separated quantile Newton lemma in
\path{../size_bounds/generalized_local_two_variable_completion.tex}
applies uniformly because all conditional density bounds and polynomial
degrees are controlled. Increase $H$ also so that $H\ge(n-k)^2$.
Each special variable now has exactly $HD_*$ equiprobable nodes;
all $k$ specified counts may therefore be chosen equal and polynomial.

On each product of pure cells the order is fixed, so the unchanged
cell moments preserve every full-order kernel. Each cell's discrete
cluster has at least $(n-k)^2$ distinct interior nodes and rational
moments through degree $n-k$. Shrink the finitely many cells around
their nodes to obtain disjoint cluster closures. The conditional lifting
theorem in
\path{../size_bounds/clustered_multi_exceptional_lifting.tex}
then converts the remaining uniforms into finite dice, retaining every
specified count $HD_*$. It provides one common finite count for those
remaining dice. Globally order the face labels and replace them by
successive integers to make all labels distinct. This proves the theorem.
\end{proof}
\end{document}
